
\section{Equivalent electrons and Hund's rules}
For two equivalent \(d\) electrons (\(d^2\)), antisymmetry permits
\[
{}^3F_{2,3,4},\quad {}^3P_{0,1,2},\quad
{}^1G_4,\quad {}^1D_2,\quad {}^1S_0.
\]
The microstate count is
\[
(3\times7)+(3\times3)+(1\times9)+(1\times5)+(1\times1)
=45=\binom{10}{2}.
\]
Hund's first rule favours maximum \(S\); second favours maximum \(L\)
within a multiplicity; third favours the smallest \(J\) for a less-than-half-full
shell. Thus the simple LS-coupled ground-state prediction for \(d^2\)
is \({}^3F_2\). Do not assume Hund's rules uniquely order every
other term: compare to experimental level data.

\section{Radial structure and nuclear-size effects}
For a central potential,
\(\psi_{nlm}(r,\theta,\phi)=R_{nl}(r)Y_{lm}(\theta,\phi)\)
and \(\int_0^\infty|R_{nl}|^2r^2\,\diff r=1\).
The alternative radial function \(u_{nl}(r)=rR_{nl}(r)\) satisfies
\(\int_0^\infty|u_{nl}|^2\,\diff r=1\).
A small change in a central potential can be estimated by first-order
perturbation theory:
\[
\Delta E_{nlm}^{(1)}
=\int\psi_{nlm}^*(\mathbf r)\,\Delta V(r)\,
\psi_{nlm}(\mathbf r)\,\diff^3r.
\]
For a spherically symmetric perturbation the angular integral is unity.
Small-\(r\) contributions matter particularly for orbitals with density
at the nucleus.

\section{E1 dipole matrix elements and radiative lifetimes}
The dipole matrix element factors into angular and radial pieces;
the radial part for \(R\)-normalisation is
\[
\int_0^\infty R_{n_fl_f}^*(r)\,r\,R_{n_il_i}(r)\,r^2\,\diff r.
\]
For specific non-degenerate initial and final states in vacuum,
the spontaneous E1 rate in SI units is
\[
A_{i\to f}=\frac{\omega_{if}^3}{3\pi\epsilon_0\hbar c^3}
\left|\langle f|\hat{\mathbf d}|i\rangle\right|^2,
\quad \hat{\mathbf d}=-e\hat{\mathbf r}.
\]
Treat sums/averages over magnetic substates consistently.
The lifetime includes \emph{all} allowed downward decay channels:
\[
\boxed{\tau_i=\left(\sum_{f:E_f<E_i} A_{i\to f}\right)^{-1}.}
\]
Atomic unit of time: \(t_{\mathrm{au}}\simeq2.418884\times10^{-17}\,\mathrm s\).
\begin{cautionbox}
Using \(u(r)\) in an integral derived for \(R(r)\) introduces incorrect
factors of \(r\). Check radial normalisation and unit conversion.
\end{cautionbox}
